\newif\ifglobal

%%%%%%%%%%%%%%%%%%%%%%%
%%% Compile to view %%%
%%%%%%%%%%%%%%%%%%%%%%%

%%%%%%%%%%%%%%%%%%%%%%%%%%%%%%%%%%%%%%%%%%%%%%%%%%%
%%% Readme:                                     %%%
%%% https://www.overleaf.com/read/bwdjvfjksvjy/ %%%
%%%%%%%%%%%%%%%%%%%%%%%%%%%%%%%%%%%%%%%%%%%%%%%%%%%

% >>> M?: marks a module setting number or important notice

% >>> M4: Set {./relative-path} to external files for this module <<<
\def \modulefiles {./46-PCNT}
% To add an external file, use:
% (Replace '---' with actual filename)
% '\input{\modulefiles/tables/---.tex}' for tables
% '\includegraphics{\modulefiles/figures/---}' for figures
% '\subfile{\modulefiles/---}' for register files


% >>> M1: This module is in global project? <<<
% 'YES' - uncomment; 'NO' - comment out
\globaltrue

\ifglobal
    \documentclass[main\_\_CN.tex]{subfiles}

\else
    % Font "FandolSong-Regular" used by ctex does not have GSUB table.
    % It causes warnings that do not affect compilation and can be ignored.
    \PassOptionsToPackage{quiet}{fontspec}

    \documentclass[openany, 10pt]{book}

    % >>> M2: In ./00-shared/preamble-trm-module, also find
    %         settings for visibility of todo notes, watermark <<<
    \usepackage{preamble-trm-module}


    % Enable Chinese typesetting
    \usepackage{ctex}

    % >>> M3: Temporary labels in Chinese
    %         you can add your own labels there <<<
    \myexternaldocument{./00-shared/temp-labels-cn}

\fi %\ifglobal


\setcounter{secnumdepth}{4}

% Define formatting of chapters in text
\titleformat{\chapter}
  [display]
  {\LARGE\bfseries}
  {\color{AllportsColor}第\, \thechapter\, 章}
  {1em}
  {\color{AllportsColor}}
\titlespacing*{\chapter}{0pt}{0pt}{20pt}

% Define headers
\usepackage{fancyhdr}
\usepackage{lipsum}
\renewcommand{\chaptermark}[1]{\markboth{第\, \thechapter\, 章 \  \ #1}{}}
\fancypagestyle{plain}{%
  \fancyhead{}%
  %\fancyhead[R]{\Navigation}
  \fancyhead[L]{%
    \nouppercase{%
      \ %
      \normalsize\leftmark
    }%
  }%
}
\pagestyle{plain}

% >>> M5: If this module needs any updates in
% './00-shared/preamble-trm-module.sty'
% (include additional package, etc.), contact your module PM
% or create the file './00-shared/changelog.tex' make the
% necessary updates yourself and document them in this file <<<

% Make text left-aligned and allow latex to break pages naturally, comment out for full justification
\raggedright
\raggedbottom


\begin{document}


% Sets language into which pre-defined bits of text are translated
\selectlanguage{chinese}

% Do not delete this block
\ifglobal
% Add the following front matter if
% this module is outside of global project
\else
    \InputIfFileExists{readme}{}{}
    \listoftodos
    {\let\clearpage\relax \listofdonetodos}
    \tableofcontents
    %\listoftables
    %\listoffigures
\fi


%%%%%%%%%%%%%%%%%%%%%%%%%
%%% TRM Module Begins %%%
%%%%%%%%%%%%%%%%%%%%%%%%%


\hypertarget{pcnt}{}
\chapter{脉冲计数控制器 (PCNT)}
\label{mod:pcnt}

\section{概述}
\label{sec:pcnt-overview}

\textbf{脉冲计数控制器 (Pulse Count Controller, PCNT)} 用于对输入脉冲计数，通过记录输入脉冲信号的上升沿或下降沿进行递增或递减计数。\chipname{} 在主系统有两个 PCNT 控制器（PCNT0 和 PCNT1）。PCNT 有四个独立的脉冲计数器，称为单元，各有一组寄存器。所有单元共用一个时钟源 APB\_CLK。

下文描述中 \regindex{n} 表示单元编号 0 \verb+~+ 3。

每个单元有两个通道（ch0 和 ch1），可以独立配置为递增或递减计数。两个通道功能相同，下文以通道 0 (ch0) 为例进行介绍。

\section{主要特性}\label{sec:pcnt-features}
PCNT 有如下特性：
\begin{itemize}
    \item 四个独立的脉冲计数器（单元），每个为 16 位有符号计数器，计数范围是 -32,768 \verb+~+ 32,767
    \item 每个单元有两个独立的通道，共用一个脉冲计数器
    \item 所有通道均有输入脉冲信号（如 sig\_ch0\_u\regindex{n}）和相应的控制信号（如 ctrl\_ch0\_u\regindex{n}）
    \item 每个单元有五个毛刺滤波器，分别过滤两个通道的输入脉冲信号 sig\_ch0\_u\regindex{n}、sig\_ch1\_u\regindex{n}，两个通道的控制信号 ctrl\_ch0\_u\regindex{n}、ctrl\_ch1\_u\regindex{n}，以及该单元的 GPIO 输入清零信号 IO\_RST\_U\regindex{n}。五个滤波器共用一套使能和阈值，按单元配置
    \item 每个通道参数如下：
        \begin{enumerate}
            \item 选择在输入脉冲信号的上升沿或下降沿计数
            \item 在控制信号为高电平或低电平时可将计数模式配置为递增模式、递减模式或停止计数
            \item 通过设置递增/递减计数步长阈值实现步长计数警报
            \item 置位清零寄存器或通过 GPIO 输入清零信号可清除脉冲计数器的值
            \item 每个计数事件观察点（预设的计数值）均可生成并记录相应的计数事件信号，该信号可触发中断
        \end{enumerate}
    \item 最大输入脉冲频率：$\frac{f_{APB\_CLK}}{2}$
\end{itemize}


\section{结构概览}
\label{sec:pcnt-arch-overview}

如图 \ref{fig:pcnt-block-diagram} \textit{\nameref{fig:pcnt-block-diagram}} 所示，每个通道有两个输入信号：
\begin{enumerate}
    \item 一个脉冲输入信号（如图 \ref{fig:pcnt-arch} \textit{\nameref{fig:pcnt-arch}} 中的 sig\_ch0\_u\regindex{n} 为单元 \regindex{n} ch0 的脉冲输入信号）
    \item 一个控制信号（如图 \ref{fig:pcnt-arch} \textit{\nameref{fig:pcnt-arch}} 中的 ctrl\_ch0\_u\regindex{n} 为单元 \regindex{n} ch0 的控制信号）
\end{enumerate}

\begin{figure}[H]
\centering
\includegraphics[width=0.4\textwidth]{\modulefiles/figures/pcnt-block-diagram-rst.eps}
\caption{PCNT 框图}
\label{fig:pcnt-block-diagram}
\end{figure}

\begin{figure}[H]
\centering
\includegraphics[width=1\textwidth]{\modulefiles/figures/pcnt_arch_updated_20250609}
\caption{PCNT 单元基本架构图}
\label{fig:pcnt-arch}
\end{figure}

如图 \ref{fig:pcnt-arch} \textit{\nameref{fig:pcnt-arch}} 所示，每个单元有通道 0 和通道 1 两个通道。每个通道处理不同的输入脉冲信号，并通过各自的 inc\_dec 模块递增或递减计数值。两个通道将计数值发送给加法器模块，该模块包含一个 16 位有符号寄存器。软件可以通过置位 \hyperref[fielddesc:PCNTCNTPAUSEUN]{PCNT\_CNT\_PAUSE\_U\regindex{n}} 暂停加法器，也可以通过置位 \hyperref[fielddesc:PCNTPULSECNTRSTUN]{PCNT\_PULSE\_CNT\_RST\_U\regindex{n}} 清零加法器。

\section{功能描述}\label{sec:pcnt-func-descr}

\subsection{计数模式}

图 \ref{fig:pcnt-arch} \textit{\nameref{fig:pcnt-arch}} 为 PCNT 单元的基本架构图。如上所述，ctrl\_ch0\_u\regindex{n} 为单元 \regindex{n} ch0 的控制信号，控制信号 ctrl\_ch0\_u\regindex{n} 为高电平或低电平时可配置不同的计数模式，在输入脉冲信号 sig\_ch0\_u\regindex{n} 的上升沿和下降沿计数。可选计数模式如下：

\begin{itemize}
\item 递增模式：通道检测到 sig\_ch0\_u\regindex{n} 的有效边沿（有效边沿软件可配）时，计数器的值 \hyperref[fielddesc:PCNTPULSECNTUN]{pulse\_cnt} 加 1。pulse\_cnt 的值达到 \hyperref[fielddesc:PCNTCNTHLIMUN]{PCNT\_CNT\_H\_LIM\_U\regindex{n}} 时被清零。如果在 pulse\_cnt 达到 \hyperref[fielddesc:PCNTCNTHLIMUN]{PCNT\_CNT\_H\_LIM\_U\regindex{n}} 前该通道的计数模式改变，则 pulse\_cnt 停止计数，新的计数模式生效。如果在 pulse\_cnt 达到 \hyperref[fielddesc:PCNTCNTHLIMUN]{PCNT\_CNT\_H\_LIM\_U\regindex{n}} 前 \hyperref[fielddesc:PCNTCNTPAUSEUN]{PCNT\_CNT\_PAUSE\_U\regindex{n}} 置 1，则 pulse\_cnt 停止计数，已配置的计数模式不变。

\item 递减模式：通道检测到 sig\_ch0\_u\regindex{n} 的有效边沿（有效边沿软件可配）时，计数器的值 pulse\_cnt 减 1。pulse\_cnt 的值达到 \hyperref[fielddesc:PCNTCNTLLIMUN]{PCNT\_CNT\_L\_LIM\_U\regindex{n}} 时被清零。如果在 pulse\_cnt 达到 \hyperref[fielddesc:PCNTCNTLLIMUN]{PCNT\_CNT\_L\_LIM\_U\regindex{n}} 前该通道的计数模式改变，则 pulse\_cnt 停止计数，新的计数模式生效。如果在 pulse\_cnt 达到 \hyperref[fielddesc:PCNTCNTLLIMUN]{PCNT\_CNT\_L\_LIM\_U\regindex{n}} 前 \hyperref[fielddesc:PCNTCNTPAUSEUN]{PCNT\_CNT\_PAUSE\_U\regindex{n}} 置 1，则 pulse\_cnt 停止计数，已配置的计数模式不变。

\item 停止计数：计数停止，计数器的值 pulse\_cnt 保持不变。
\end{itemize}


表 \ref{tab:pcnt-positive-low} \textit{\nameref{tab:pcnt-positive-low}} 至表 \ref{tab:pcnt-negative-high} \textit{\nameref{tab:pcnt-negative-high}} 说明了如何配置通道 0 的计数模式。

\begin{table}[H]
\centering
\caption{控制信号为低电平时输入脉冲信号上升沿的计数模式}
\label{tab:pcnt-positive-low}
\begin{tabular}{|l|l|l|}
\hline
\rowcolor{lightgray}
\textbf{\hyperref[fielddesc:PCNTCH0POSMODEUN]{PCNT\_CH0\_POS\_MODE\_U\regindex{n}}} & \textbf{\hyperref[fielddesc:PCNTCH0LCTRLMODEUN]{PCNT\_CH0\_LCTRL\_MODE\_U\regindex{n}}} & \textbf{计数模式} \\ \hline
\multirow{3}{*}{1}             & 0                                     & 递增模式    \\ \cline{2-3}
                               & 1                                     & 递减模式    \\ \cline{2-3}
                               & 其他                                & 停止计数      \\ \hline
\multirow{3}{*}{2}             & 0                                     & 递减模式    \\ \cline{2-3}
                               & 1                                     & 递增模式    \\ \cline{2-3}
                               & 其他                                & 停止计数      \\ \hline
其他                         & N/A                                   & 停止计数      \\ \hline
\end{tabular}
\end{table}

\begin{table}[H]
\centering
\caption{控制信号为高电平时输入脉冲信号上升沿的计数模式}
\label{tab:pcnt-positive-high}
\begin{tabular}{|l|l|l|}
\hline
\rowcolor{lightgray}
\textbf{\hyperref[fielddesc:PCNTCH0POSMODEUN]{PCNT\_CH0\_POS\_MODE\_U\regindex{n}}} & \textbf{\hyperref[fielddesc:PCNTCH0HCTRLMODEUN]{PCNT\_CH0\_HCTRL\_MODE\_U\regindex{n}}} & \textbf{计数模式} \\ \hline
\multirow{3}{*}{1}             & 0                                     & 递增模式    \\ \cline{2-3}
                               & 1                                     & 递减模式    \\ \cline{2-3}
                               & 其他                                & 停止计数      \\ \hline
\multirow{3}{*}{2}             & 0                                     & 递减模式    \\ \cline{2-3}
                               & 1                                     & 递增模式    \\ \cline{2-3}
                               & 其他                                & 停止计数      \\ \hline
其他                         & N/A                                   & 停止计数      \\ \hline
\end{tabular}
\end{table}

\begin{table}[H]
\centering
\caption{控制信号为低电平时输入脉冲信号下降沿的计数模式}
\label{tab:pcnt-negative-low}
\begin{tabular}{|l|l|l|}
\hline
\rowcolor{lightgray}
\textbf{\hyperref[fielddesc:PCNTCH0NEGMODEUN]{PCNT\_CH0\_NEG\_MODE\_U\regindex{n}}} & \textbf{\hyperref[fielddesc:PCNTCH0LCTRLMODEUN]{PCNT\_CH0\_LCTRL\_MODE\_U\regindex{n}}} & \textbf{计数模式} \\ \hline
\multirow{3}{*}{1}             & 0                                     & 递增模式    \\ \cline{2-3}
                               & 1                                     & 递减模式    \\ \cline{2-3}
                               & 其他                                & 停止计数      \\ \hline
\multirow{3}{*}{2}             & 0                                     & 递减模式    \\ \cline{2-3}
                               & 1                                     & 递增模式    \\ \cline{2-3}
                               & 其他                                & 停止计数      \\ \hline
其他                         & N/A                                   & 停止计数      \\ \hline
\end{tabular}
\end{table}

\begin{table}[H]
\centering
\caption{控制信号为高电平时输入脉冲信号下降沿的计数模式}
\label{tab:pcnt-negative-high}
\begin{tabular}{|l|l|l|}
\hline
\rowcolor{lightgray}
\textbf{\hyperref[fielddesc:PCNTCH0NEGMODEUN]{PCNT\_CH0\_NEG\_MODE\_U\regindex{n}}} & \textbf{\hyperref[fielddesc:PCNTCH0HCTRLMODEUN]{PCNT\_CH0\_HCTRL\_MODE\_U\regindex{n}}} & \textbf{计数模式} \\ \hline
\multirow{3}{*}{1}             & 0                                     & 递增模式    \\ \cline{2-3}
                               & 1                                     & 递减模式    \\ \cline{2-3}
                               & 其他                                & 停止计数      \\ \hline
\multirow{3}{*}{2}             & 0                                     & 递减模式    \\ \cline{2-3}
                               & 1                                     & 递增模式    \\ \cline{2-3}
                               & 其他                                & 停止计数      \\ \hline
其他                         & N/A                                   & 停止计数      \\ \hline
\end{tabular}
\end{table}

\subsection{脉冲毛刺滤波器}

每个单元有五个毛刺滤波器，分别用于滤除两个通道的输入脉冲信号 (sig\_ch0\_u\regindex{n}、sig\_ch1\_u\regindex{n})、两个通道的控制信号 (ctrl\_ch0\_u\regindex{n}、ctrl\_ch1\_u\regindex{n}) 以及该单元的 GPIO 输入清零信号 (IO\_RST\_U\regindex{n}) 的毛刺。五个滤波器共享同一组使能配置和阈值配置，置位 \hyperref[fielddesc:PCNTFILTERENUN]{PCNT\_FILTER\_EN\_U\regindex{n}} 同时使能五个滤波器，滤除脉冲宽度小于 \hyperref[fielddesc:PCNTFILTERTHRESUN]{PCNT\_FILTER\_THRES\_U\regindex{n}} 个 APB 时钟周期的线路毛刺。

\subsection{配置更新}

\begin{itemize}

\item 在 PCNT 模块工作过程中，软件如果重新配置 \hyperref[fielddesc:PCNTCNTHLIMUN]{PCNT\_CNT\_H\_LIM\_U\regindex{n}} 和/或 \hyperref[fielddesc:PCNTCNTLLIMUN]{PCNT\_CNT\_L\_LIM\_U\regindex{n}}，则新的配置值会在 pulse\_cnt 达到最大计数值和最小计数值的任意一个后再生效。

\item 软件如果重新配置 \hyperref[fielddesc:PCNTCNTTHRES0UN]{PCNT\_CNT\_THRES0\_U\regindex{n}}、\hyperref[fielddesc:PCNTCNTTHRES1UN]{PCNT\_CNT\_THRES1\_U\regindex{n}}、\hyperref[fielddesc:PCNTCNTHSTEPUN]{PCNT\_CNT\_H\_STEP\_U\regindex{n}} 和/或 \hyperref[fielddesc:PCNTCNTLSTEPUN]{PCNT\_CNT\_L\_STEP\_U\regindex{n}}，则新的配置值会立即生效。
\end{itemize}

\section{中断和计数事件观察点}

PCNT 可以设置七个计数事件观察点。这些观察点共享一个中断。每个观察点都有独立的中断使能信号，用于开启或屏蔽中断。

\begin{itemize}
    \item 最大计数值：当 pulse\_cnt 大于等于 \hyperref[fielddesc:PCNTCNTHLIMUN]{PCNT\_CNT\_H\_LIM\_U\regindex{n}} 时，产生上限中断，同时将 \hyperref[fielddesc:PCNTCNTTHRHLIMLATUN]{PCNT\_CNT\_THR\_H\_LIM\_LAT\_U\regindex{n}} 置高。
    \item 最小计数值：当 pulse\_cnt 小于等于 \hyperref[fielddesc:PCNTCNTLLIMUN]{PCNT\_CNT\_L\_LIM\_U\regindex{n}} 时，产生下限中断，同时将 \hyperref[fielddesc:PCNTCNTTHRLLIMLATUN]{PCNT\_CNT\_THR\_L\_LIM\_LAT\_U\regindex{n}} 置高。
    \item 两个中间阈值：当 pulse\_cnt 等于 \hyperref[fielddesc:PCNTCNTTHRES0UN]{PCNT\_CNT\_THRES0\_U\regindex{n}} 或者 \hyperref[fielddesc:PCNTCNTTHRES1UN]{PCNT\_CNT\_THRES1\_U\regindex{n}} 时，产生中断，同时将 \hyperref[fielddesc:PCNTCNTTHRTHRES0LATUN]{PCNT\_CNT\_THR\_THRES0\_LAT\_U\regindex{n}} 或 \hyperref[fielddesc:PCNTCNTTHRTHRES1LATUN]{PCNT\_CNT\_THR\_THRES1\_LAT\_U\regindex{n}} 置高。
    \item 零：当 pulse\_cnt 等于 0 时，产生中断，同时 \hyperref[fielddesc:PCNTCNTTHRZEROLATUN]{PCNT\_CNT\_THR\_ZERO\_LAT\_U\regindex{n}} 有效。
    \item 递增计数步长阈值：如果 \hyperref[fielddesc:PCNTDALTACHANGEENUN]{PCNT\_DALTA\_CHANGE\_EN\_U\regindex{n}} 被设置为高电平，在计数递增过程中 pulse\_cnt 每增加 \hyperref[fielddesc:PCNTCNTHSTEPUN]{PCNT\_CNT\_H\_STEP\_U\regindex{n}} 时，将触发递增计数步长中断，并将 \hyperref[fielddesc:PCNTCNTTHRHSTEPLATUN]{PCNT\_CNT\_THR\_H\_STEP\_LAT\_U\regindex{n}} 置为高电平。

    例如如果 \hyperref[fielddesc:PCNTCNTHSTEPUN]{PCNT\_CNT\_H\_STEP\_U\regindex{n}} 设置为 100，那么当 pulse\_cnt 脉冲数达到 100、200、300 等整数倍值，将分别触发中断一次。此类中断可用于周期性地检测输入脉冲数量是否达到预定倍数，便于实现定量采样、数据统计或定期事件响应等功能。

    \item 递减计数步长阈值：如果 \hyperref[fielddesc:PCNTDALTACHANGEENUN]{PCNT\_DALTA\_CHANGE\_EN\_U\regindex{n}} 被设置为高电平，在计数递减过程中 pulse\_cnt 每减少 \hyperref[fielddesc:PCNTCNTLSTEPUN]{PCNT\_CNT\_L\_STEP\_U\regindex{n}} 时，将触发递减计数步长中断，并将 \hyperref[fielddesc:PCNTCNTTHRLSTEPLATUN]{PCNT\_CNT\_THR\_L\_STEP\_LAT\_U\regindex{n}} 置为高电平。

    例如，如果 \hyperref[fielddesc:PCNTCNTLSTEPUN]{PCNT\_CNT\_L\_STEP\_U\regindex{n}} 设置为 100，那么当 pulse\_cnt 脉冲数递减至 --100、-200、-300 等整数倍时，将分别触发中断一次。此类中断适用于周期性检测计数递减是否达到预定倍数，便于实现定量采样、数据统计或定期事件响应等功能。
\end{itemize}


\section{应用实例}

每个单元的通道 0 和通道 1 可配置为独立工作或一起工作。下文详细说明了通道 0 独自递增计数、通道 0 独自递减计数和两个通道一起递增计数的应用实例。本节中未详述的通道工作模式（如通道 1 独自递减或递增、双通道一增一减），可参考这三种模式。

\subsection{通道 0 独自递增计数} \label{sec:pcnt-func-descr-one-channel-inc}

\begin{figure}[H]
\centering
\fbox{\includegraphics[width=1\textwidth]{\modulefiles/figures/pulse_inc}}
\caption{通道 0 递增计数图}
\label{fig:pcnt-inc}
\end{figure}

图 \ref{fig:pcnt-inc} \textit{\nameref{fig:pcnt-inc}} 为通道 0 在 sig\_ch0\_u\regindex{n} 上升沿独立递增计数的示意图，此时通道 1 关闭（请参阅表 \ref{tab:pcnt-positive-low} \textit{\nameref{tab:pcnt-positive-low}} 至表 \ref{tab:pcnt-negative-high} \textit{\nameref{tab:pcnt-negative-high}}，通道 1 使用对应的 CH1 字段关闭）。通道 0 的配置如下所示。

\begin{itemize}
    \item \hyperref[fielddesc:PCNTCH0LCTRLMODEUN]{PCNT\_CH0\_LCTRL\_MODE\_U\regindex{n}}=0：当 ctrl\_ch0\_u\regindex{n} 为低电平时，递增计数。
    \item \hyperref[fielddesc:PCNTCH0HCTRLMODEUN]{PCNT\_CH0\_HCTRL\_MODE\_U\regindex{n}}=2：当 ctrl\_ch0\_u\regindex{n} 为高电平时，停止计数。
    \item \hyperref[fielddesc:PCNTCH0POSMODEUN]{PCNT\_CH0\_POS\_MODE\_U\regindex{n}}=1：在 sig\_ch0\_u\regindex{n} 的上升沿递增计数。
    \item \hyperref[fielddesc:PCNTCH0NEGMODEUN]{PCNT\_CH0\_NEG\_MODE\_U\regindex{n}}=0：在 sig\_ch0\_u\regindex{n} 的下降沿不计数。
    \item \hyperref[fielddesc:PCNTCNTHLIMUN]{PCNT\_CNT\_H\_LIM\_U\regindex{n}}=5：pulse\_cnt 的值递增至 \hyperref[fielddesc:PCNTCNTHLIMUN]{PCNT\_CNT\_H\_LIM\_U\regindex{n}} 时被清零。
\end{itemize}

\subsection{通道 0 独自递减计数} \label{sec:pcnt-one-channel-dec}

\begin{figure}[H]
\centering
\fbox{\includegraphics[width=1\textwidth]{\modulefiles/figures/pulse_dec}}
\caption{通道 0 递减计数图}
\label{fig:pcnt-dec}
\end{figure}

图 \ref{fig:pcnt-dec} \textit{\nameref{fig:pcnt-dec}} 为通道 0 在 sig\_ch0\_u\regindex{n} 上升沿独立递减计数的示意图，此时通道 1 关闭。此时通道 0 的配置与图 \ref{fig:pcnt-inc} \textit{\nameref{fig:pcnt-inc}} 相比有如下区别：

\begin{itemize}
    \item \hyperref[fielddesc:PCNTCH0POSMODEUN]{PCNT\_CH0\_POS\_MODE\_U\regindex{n}}=2：即在 sig\_ch0\_u\regindex{n} 的上升沿递减计数。
    \item \hyperref[fielddesc:PCNTCNTLLIMUN]{PCNT\_CNT\_L\_LIM\_U\regindex{n}}=-5：pulse\_cnt 的值递减到 \hyperref[fielddesc:PCNTCNTLLIMUN]{PCNT\_CNT\_L\_LIM\_U\regindex{n}} 时被清零。
\end{itemize}

\subsection{通道 0 和通道 1 同时递增计数}\label{sec:pcnt-two-channels-inc}

\begin{figure}[H]
\centering
\fbox{\includegraphics[width=1\textwidth]{\modulefiles/figures/pulse_add_2}}
\caption{双通道递增计数图}
\label{fig:pcnt-two-inc}
\end{figure}

图 \ref{fig:pcnt-two-inc} \textit{\nameref{fig:pcnt-two-inc}} 为通道 0 和通道 1 分别在 sig\_ch0\_u\regindex{n} 和 sig\_ch1\_u\regindex{n} 上升沿同时递增计数的示意图。如图 \ref{fig:pcnt-two-inc} \textit{\nameref{fig:pcnt-two-inc}} 所示，控制信号 ctrl\_ch0\_u\regindex{n} 与 ctrl\_ch1\_u\regindex{n} 的波形一致，输入脉冲信号 sig\_ch0\_u\regindex{n} 和 sig\_ch1\_u\regindex{n} 波形一致。具体配置如下：

\begin{itemize}
    \item 通道 0：
    \begin{itemize}
        \item \hyperref[fielddesc:PCNTCH0LCTRLMODEUN]{PCNT\_CH0\_LCTRL\_MODE\_U\regindex{n}}=0：当 ctrl\_ch0\_u\regindex{n} 为低电平时，递增计数。
        \item \hyperref[fielddesc:PCNTCH0HCTRLMODEUN]{PCNT\_CH0\_HCTRL\_MODE\_U\regindex{n}}=2：当 ctrl\_ch0\_u\regindex{n} 为高电平时，停止计数。
        \item \hyperref[fielddesc:PCNTCH0POSMODEUN]{PCNT\_CH0\_POS\_MODE\_U\regindex{n}}=1：在 sig\_ch0\_u\regindex{n} 的上升沿递增计数。
        \item \hyperref[fielddesc:PCNTCH0NEGMODEUN]{PCNT\_CH0\_NEG\_MODE\_U\regindex{n}}=0：在 sig\_ch0\_u\regindex{n} 的下降沿不计数。

    \end{itemize}

    \item 通道 1：
    \begin{itemize}
        \item \hyperref[fielddesc:PCNTCH1LCTRLMODEUN]{PCNT\_CH1\_LCTRL\_MODE\_U\regindex{n}}=0：当 ctrl\_ch1\_u\regindex{n} 为低电平时，递增计数。

        \item \hyperref[fielddesc:PCNTCH1HCTRLMODEUN]{PCNT\_CH1\_HCTRL\_MODE\_U\regindex{n}}=2：当 ctrl\_ch1\_u\regindex{n} 为高电平时，停止计数。

        \item \hyperref[fielddesc:PCNTCH1POSMODEUN]{PCNT\_CH1\_POS\_MODE\_U\regindex{n}}=1：在 sig\_ch1\_u\regindex{n} 的上升沿递增计数。

        \item \hyperref[fielddesc:PCNTCH1NEGMODEUN]{PCNT\_CH1\_NEG\_MODE\_U\regindex{n}}=0：在 sig\_ch1\_u\regindex{n} 的下降沿不计数。

    \end{itemize}
    \item \hyperref[fielddesc:PCNTCNTHLIMUN]{PCNT\_CNT\_H\_LIM\_U\regindex{n}}=10：pulse\_cnt 递增至 \hyperref[fielddesc:PCNTCNTHLIMUN]{PCNT\_CNT\_H\_LIM\_U\regindex{n}} 时被清零。

\end{itemize}


% >>> If your module has no registers, delete sample text
%       for Sections Register Summary and Registers below <<<

%%%%%%%%%%%%%%%%%%%%%%%%
%%% Register Summary %%%
%%%%%%%%%%%%%%%%%%%%%%%%

% >>> Update the sample text below and insert
%     the info on register summary into the subfile <<<

\hypertarget{pcnt-reg-summ}{}
\section{寄存器列表}
\label{sec:pcnt-reg-summ}

本小节的所有地址均为相对于 PCNT 基地址的地址偏移量（相对地址），具体基地址请见章节 \ref{mod:sysmem} \textit{\nameref{mod:sysmem}} 中的表 \ref{tab:sysmem-base-address} \textit{\nameref{tab:sysmem-base-address}}。

请查看章节 \hyperref[glossary-access-types]{\textit{寄存器的访问类型}}，了解“访问”列缩写的含义。

\subfile{\modulefiles/46-PCNT-reg-summ__CN}


%%%%%%%%%%%%%%%%%
%%% Registers %%%
%%%%%%%%%%%%%%%%%

% >>> Update the sample text below and insert
%      the info on registers into the subfile <<<

\newpage
\section{寄存器}

本小节的所有地址均为相对于 PCNT 基地址的地址偏移量（相对地址），具体基地址请见章节 \ref{mod:sysmem} \textit{\nameref{mod:sysmem}} 中的表 \ref{tab:sysmem-base-address} \textit{\nameref{tab:sysmem-base-address}}。

\subfile{\modulefiles/46-PCNT-reg__CN}


\end{document}
